\chapter{The Universal Tensor as a Residual Type}
\label{chap:universal-tensor-residual-type}

The universal tensor is not a mysterious bulk object to the compiler. It is the gathered type of all
contractions the construction is willing to make explicit. Each index, symmetry, and dual pairing becomes
a field the checker may demand. Each conservation law becomes a compatibility obligation. Each admissible
state determines a residual by feeding that tensor into a test. The tensor is universal not because the
compiler sees infinity at once, but because every finite reading it accepts is routed through the same
structured source.

The analytic notation can be written as a map
\[
  \mathcal U(u) \in V',
\]
where \(u\) is the state and \(V'\) is the dual space of tests. A forcing term or boundary load is another
element \(\ell \in V'\). The residual is then
\[
  R(u)(v) = \langle \mathcal U(u), v\rangle - \ell(v),
  \qquad v\in V .
\]
For analysis, this says the equation is solved when the residual vanishes on the admissible tests. For the
compiler, it says something sharper: every test \(v\) offered to the machine generates an obligation
\(R(u)(v)=0\), and that obligation has a type, a cost, and a proof state. The universal tensor is therefore
the record from which the obligations are projected.

\section{Rows, contractions, and obligations}
\label{se:rows-contractions-obligations}

A tensor contraction is a row in the compiler's ledger. It names the pieces being paired, the variance of
each slot, the metric or density needed to compare them, and the scalar output that must be checked. Nothing
about that is metaphorical. A proof assistant cannot consume a hand wave about a tensor. It can consume a
finite structure whose projections state the laws. Once those projections exist, each row is as ordinary as
any other elaboration problem: synthesize the instances, normalize the expression, expose the residual, and
decide whether the demanded equality follows.

The universal tensor becomes important precisely because it prevents rows from being invented ad hoc. A
single test row is not allowed to carry a private geometry. It must obtain its pairing, measure, boundary
behavior, and symmetry from the same source as every other row. That is the compiler's version of
coordinate-free seriousness. The expression may be rendered in coordinates for computation, but the
contracts that license those coordinates live above the rendering.

\section{Sobolev norm as the residual gauge}
\label{se:sobolev-residual-gauge}

The Sobolev norm supplies the residual gauge. Pointwise residuals are too sharp for the weak problem; they
ask for more regularity than the state may honestly have. The dual Sobolev norm instead asks how large the
residual is as a functional on admissible tests:
\[
  \|R(u)\|_{H^{-1}}
    =
  \sup_{v\ne 0}
    \frac{|R(u)(v)|}{\|v\|_{H^1}} .
\]
The compiler cannot take that supremum by magic. What it can do is preserve the type of the norm and
record finite lower or exact readings against chosen tests. When a finite Galerkin family is in force, the
same formula collapses to a matrix computation over the selected basis, and the residual gauge becomes a
finite object the checker can inspect.

This is why the norm belongs in the compiler discussion. A norm is not only a number placed after the
analysis is complete. It is the rule that tells the machine when two residuals may be compared. Without the
Sobolev gauge, a small pointwise-looking expression may be meaningless, and a distributional residual may
look larger than it is. With the gauge named, the compiler has a stable unit for activity: it can say which
residual is being measured, against which tests, in which topology, and at what cost.

\section{Universality without hidden choice}
\label{se:universality-without-hidden-choice}

Universality is dangerous when it conceals selection. The compiler has already learned to distrust hidden
choice: a result that depends on an unnamed witness is not the same result as one that carries its witness
in the term. The universal tensor avoids that danger only if its finite readings are explicit. A basis is
named. A quadrature rule is named, or absent. A boundary trace is named. A symmetry reduction is named. The
resulting residual is no longer an informal impression of the continuum equation; it is a typed projection
from the universal object into a finite ledger entry.

The compiler-facing tensor therefore has two jobs. It gathers the invariant structure so that every row
uses the same law, and it exposes enough finite projections for exact computation. The first job protects
meaning. The second protects checkability. Neither is subordinate to the other. A compiler that can check
finite rows without knowing what they are rows of has only arithmetic. A compiler that knows the invariant
object but cannot project it into rows has only syntax. The universal tensor is the hinge between the two.
